\documentclass[12pt]{article}
\usepackage{amsmath, amssymb}
\usepackage{enumitem}
\usepackage{geometry}
\geometry{margin=1in}

\begin{document}

\section*{Lecture 5: Bayes’ Theorem, Random Variables, and Probability Mass Function}

\hrulefill

\subsection*{1. Introduction and Scope of the Lecture}

This lecture covers three major topics:

\begin{enumerate}
    \item Bayes’ Theorem
    \item Random Variables
    \item Probability Mass Function (PMF)
\end{enumerate}

The focus is on understanding conditional probability through Bayes’ theorem, motivating the concept of random variables, and formally defining probability mass functions for discrete random variables.

\hrulefill

\section*{Part I: Bayes’ Theorem}

\subsection*{2. Bayes’ Theorem as a Weighted Average of Conditional Probabilities}

\subsubsection*{Definition and Decomposition of an Event}

Let \textbf{A} and \textbf{B} be two events. Any outcome that belongs to event \textbf{A} must satisfy one of the following:

\begin{itemize}
    \item It lies in both \textbf{A} and \textbf{B} (event \textbf{AB}), or
    \item It lies in \textbf{A} but not in \textbf{B} (event \textbf{AB}$^{c}$).
\end{itemize}

Thus, event \textbf{A} can be written as:

\[
A = AB \cup AB^{c}
\]

\subsubsection*{Mutual Exclusivity}

The events \textbf{AB} and \textbf{AB}$^{c}$ are mutually exclusive because an outcome cannot simultaneously belong to \textbf{B} and \textbf{B}$^{c}$.

\subsubsection*{Application of Axiom 3}

By Axiom 3 of probability (additivity for mutually exclusive events):

\[
\Pr(A) = \Pr(AB) + \Pr(AB^{c})
\]

\subsubsection*{Expressing Joint Probabilities Using Conditional Probability}

Using the definition of conditional probability:

\[
\Pr(AB) = \Pr(A \mid B)\Pr(B)
\]

\[
\Pr(AB^{c}) = \Pr(A \mid B^{c})\Pr(B^{c})
\]

Since $\Pr(B^{c}) = 1 - \Pr(B)$, we obtain:

\[
\Pr(A) = \Pr(A \mid B)\Pr(B) + \Pr(A \mid B^{c})[1 - \Pr(B)]
\]

\subsubsection*{Interpretation}

The probability of event \textbf{A} is a \textbf{weighted average} of the conditional probabilities $\Pr(A \mid B)$ and $\Pr(A \mid B^{c})$, where the weights are the probabilities of \textbf{B} and \textbf{B}$^{c}$, respectively.

\hrulefill

\subsection*{3. Learning by Example: Insurance Policy Example (Example 3.1)}

\subsubsection*{Given Information}

\begin{itemize}
    \item Population divided into two classes:
    \begin{itemize}
        \item Accident-prone
        \item Not accident-prone
    \end{itemize}
    \item Probability of accident within one year:
    \begin{itemize}
        \item Accident-prone: 0.4
        \item Not accident-prone: 0.2
    \end{itemize}
    \item Probability that a randomly chosen person is accident-prone: 0.3
\end{itemize}

\subsubsection*{Part 1: Probability That a New Policyholder Has an Accident}

Let:

\begin{itemize}
    \item $A_1$ = event that the policyholder has an accident within one year
    \item $A$ = event that the policyholder is accident-prone
\end{itemize}

Using conditioning on whether the policyholder is accident-prone:

\[
\Pr(A_1) = \Pr(A_1 \mid A)\Pr(A) + \Pr(A_1 \mid A^{c})\Pr(A^{c})
\]

Substituting the given values:

\[
\Pr(A_1) = (0.4)(0.3) + (0.2)(0.7)
\]

\[
\Pr(A_1) = 0.12 + 0.14 = 0.26
\]

Thus, the probability that a new policyholder has an accident within one year is \textbf{0.26}.

\hrulefill

\subsubsection*{Part 2: Probability That the Policyholder Is Accident-Prone Given an Accident Occurred}

Now, it is given that the policyholder \textbf{did} have an accident within one year.

We want to find:

\[
\Pr(A \mid A_1)
\]

Using the definition of conditional probability:

\[
\Pr(A \mid A_1) = \frac{\Pr(A \cap A_1)}{\Pr(A_1)}
\]

Using joint probability:

\[
\Pr(A \cap A_1) = \Pr(A)\Pr(A_1 \mid A) = (0.3)(0.4)
\]

Thus:

\[
\Pr(A \mid A_1) = \frac{0.3 \times 0.4}{0.26} = \frac{6}{13}
\]

\hrulefill

\subsection*{4. Formal Introduction: Law of Total Probability and Bayes’ Formula}

\subsubsection*{Law of Total Probability}

Let $B_1, B_2, \ldots, B_n$ be mutually exclusive and exhaustive events such that one of them must occur.

Then for any event $A$:

\[
\Pr(A) = \sum \Pr(A \cap B_i)
\]

Using conditional probability:

\[
\Pr(A \cap B_i) = \Pr(A \mid B_i)\Pr(B_i)
\]

Hence:

\[
\Pr(A) = \sum \Pr(A \mid B_i)\Pr(B_i)
\]

This is known as the \textbf{Law of Total Probability}.

\hrulefill

\subsubsection*{Bayes’ Formula (Proposition 3.1)}

Starting from:

\[
\Pr(A \cap B_i) = \Pr(B_i \mid A)\Pr(A)
\]

And substituting into the Law of Total Probability, we obtain:

\[
\Pr(B_i \mid A) =
\frac{\Pr(A \mid B_i)\Pr(B_i)}
{\sum \Pr(A \mid B_j)\Pr(B_j)}
\]

\subsubsection*{Terminology}

\begin{itemize}
    \item \textbf{$\Pr(B_i)$}: \textit{A priori probability} (formed from prior assumptions or models)
    \item \textbf{$\Pr(B_i \mid A)$}: \textit{A posteriori probability} (updated after observing event A)
\end{itemize}

\hrulefill

\subsection*{5. Learning by Example: Card Experiment (Example 3.2)}

\subsubsection*{Description of the Experiment}

Three cards:

\begin{itemize}
    \item One card: Red--Red (RR)
    \item One card: Black--Black (BB)
    \item One card: Red--Black (RB)
\end{itemize}

One card is randomly selected and placed on the ground. The upper side is observed to be red.

Question: What is the probability that the other side is black?

\hrulefill

\subsubsection*{Event Definitions}

\begin{itemize}
    \item RR: card is red--red
    \item BB: card is black--black
    \item RB: card is red--black
    \item R: upper side is red
\end{itemize}

Each card is equally likely:
\[
\Pr(RR) = \Pr(BB) = \Pr(RB) = \frac{1}{3}
\]

\hrulefill

\subsubsection*{Applying Bayes’ Formula}

\[
\Pr(RB \mid R) =
\frac{\Pr(R \mid RB)\Pr(RB)}
{\Pr(R \mid RR)\Pr(RR) + \Pr(R \mid RB)\Pr(RB) + \Pr(R \mid BB)\Pr(BB)}
\]

Substitute values:

\begin{itemize}
    \item $\Pr(R \mid RR) = 1$
    \item $\Pr(R \mid RB) = \frac{1}{2}$
    \item $\Pr(R \mid BB) = 0$
\end{itemize}

Thus:

\[
\Pr(RB \mid R) =
\frac{\frac{1}{2} \times \frac{1}{3}}
{\left(1 \times \frac{1}{3}\right) + \left(\frac{1}{2} \times \frac{1}{3}\right) + \left(0 \times \frac{1}{3}\right)}
\]

\[
\Pr(RB \mid R) = \frac{1}{6} \div \frac{1}{2} = \frac{1}{3}
\]

\hrulefill

\subsubsection*{Explanation of the Common Mistake}

It is incorrect to assume that the probabilities of RR and RB are equal after observing a red side.

Considering individual sides, there are six equally likely outcomes:
$R_1, R_2, B_1, B_2, R_3, B_3$.

Given a red side is observed, the possible outcomes are $R_1, R_2,$ and $R_3$.

Only $R_3$ corresponds to the red--black card, so the probability is $1/3$.

\hrulefill

\section*{Part II: Random Variables}

\subsection*{6. Motivation and Concept}

In many experiments, interest lies not in the exact outcome, but in a function of the outcome.

Examples:

\begin{itemize}
    \item In dice tossing, interest may lie in the \textbf{sum} of the dice.
    \item In coin tossing, interest may lie in the \textbf{number of heads}.
\end{itemize}

A \textbf{random variable} is a real-valued function defined on the sample space.

\begin{itemize}
    \item Its value is determined by the outcome of the experiment.
    \item Probabilities are assigned to its possible values.
\end{itemize}

The distribution of a random variable can be visualized using a bar diagram:

\begin{itemize}
    \item X-axis: values of the random variable
    \item Height of each bar: $\Pr[X = a]$
\end{itemize}

\hrulefill

\subsection*{7. Example: Tossing Three Fair Coins}

Let the experiment consist of tossing three fair coins.

Let $Y$ denote the number of heads observed.

Possible values of $Y$: 0, 1, 2, 3

Corresponding probabilities:

\begin{itemize}
    \item $\Pr(Y = 0) = \Pr(TTT) = \frac{1}{8}$
    \item $\Pr(Y = 1) = \Pr(TTH, THT, HTT) = \frac{3}{8}$
    \item $\Pr(Y = 2) = \Pr(THH, HTH, HHT) = \frac{3}{8}$
    \item $\Pr(Y = 3) = \Pr(HHH) = \frac{1}{8}$
\end{itemize}

Since $Y$ must take one of these values, the sum of probabilities equals 1.

\hrulefill

\section*{Part III: Probability Mass Function (PMF)}

\subsection*{8. Definition of a Discrete Random Variable}

A random variable that can take on at most a \textbf{countable number} of possible values is called a \textbf{discrete random variable}.

\hrulefill

\subsection*{9. Probability Mass Function}

Let $X$ be a discrete random variable with range:

\[
R_X = \{x_1, x_2, x_3, \ldots\}
\]

The function:

\[
p(x_k) = \Pr(X = x_k)
\]

is called the \textbf{Probability Mass Function (PMF)} of $X$.

Since $X$ must take one of the values in its range:

\[
\sum p(x_k) = 1
\]

\hrulefill

\subsection*{10. Example: PMF Definition Problem}

Given the PMF:

\[
p(i) = c \frac{\lambda^i}{i!}, \quad i = 0, 1, 2, \ldots, \text{ where } \lambda > 0
\]

Using the PMF definition and the fact that probabilities sum to 1, required probabilities such as:

\begin{itemize}
    \item $\Pr(X = 0)$
    \item $\Pr(X > 2)$
\end{itemize}

are obtained by substituting the appropriate values of $i$ into the given PMF and summing over the required range.

\hrulefill

\subsection*{End of Lecture 5 Scribe}

\end{document}
